\section{Einleitung}
Bei der Robotersimulation geht es darum das Verhalten eines Roboters in einer Softwareumgebung virtuell nachzubilden und
somit das Testen und Programmieren des Roboters außerhalb der Produktivumgebung zu ermöglichen.
Dabei werden häufig grafische Softwaresysteme verwendet, die die Umgebung und Bewegung des Roboters visuell darstellen.
Gerade in der Offline-Programmierung spielt die Robotersimulation eine zentrale Rolle,
um die textuelle Programmierung in Form von CAD-basierter Programmierung zu ergänzen.
Auf diese Weise lässt sich das zeitaufwendige Teachen in der Produktionsumgebung deutlich reduzieren \doublecite[113, 122]{weber-Industrieroboter}.
 
Im Rahmen dieser Arbeit wird ein Framework zur Robotersimulation entwickelt,
das innerhalb der JupyterLab-Umgebung genutzt werden kann und in erster Linie für den Einsatz in der Lehre vorgesehen ist.
Das Framework wird anschließend auf einem JupyterHub-Server als Kernel eingerichtet,
sodass es Stdudierenden zur Verfügung steht.

Die vorliegende Arbeit baut auf den Ergebnissen einer vorausgegangenen Projektarbeit auf,
in der verschiedene Grafikbibliotheken evaluiert sowie grundlegende Funktionalitäten wie die Erstellung einer 3D-Umgebung,
Transformationen und Animationen von Körpern umgesetzt wurden. Auch die Simulation der Bewegung eines radgetriebenen Roboters wurde im Rahmen dieser Projektarbeit realisiert.
Aufbauend auf diesen Ergebnissen wird das entwickelte Framework in dieser Arbeit um die Steuerung und Animation beliebiger Industrieroboter mit serieller Kinematik erweitert.

Ziel dieser Arbeit ist die Erweiterung bestehender Jupyter-Notebooks im Lehrbetrieb der Fachhochschule Dortmund um visuelle und interaktive Komponenten, insbesondere im technischen Modulbereich wie der Robotik.
Durch die Integration anschaulicher Visualisierungen und direkter Benutzerinteraktionen soll das Verständnis kinematischer Zusammenhänge gefördert und die Lehre praxisnäher gestaltet werden.
